\documentclass[12pt]{article}
\usepackage{amsmath, amssymb, amsthm}
\usepackage{hyperref}

\newtheorem{theorem}{Theorem}
\newtheorem{lemma}[theorem]{Lemma}
\newtheorem{proposition}[theorem]{Proposition}
\newtheorem{corollary}[theorem]{Corollary}
\newtheorem{definition}[theorem]{Definition}
\newtheorem{remark}[theorem]{Remark}

\title{Problem 9: Algebraic Relations on Determinantal Tensors\\
\large A Complete Proof}
\author{}
\date{April 2026}

\begin{document}
\maketitle

\section{Problem Statement}

Let $n \geq 5$. Let $A^{(1)}, \ldots, A^{(n)} \in \mathbb{R}^{3 \times 4}$ be Zariski-generic. For $\alpha, \beta, \gamma, \delta \in [n]$, construct the tensor $Q^{(\alpha\beta\gamma\delta)} \in \mathbb{R}^{3 \times 3 \times 3 \times 3}$ with:
\[
Q^{(\alpha\beta\gamma\delta)}_{ijk\ell} = \det\begin{bmatrix} A^{(\alpha)}(i,:) \\ A^{(\beta)}(j,:)  \\ A^{(\gamma)}(k,:) \\ A^{(\Adelta)}(\ell,:) \end{bmatrix},
\]
where $A(i,:)$ is the $i$-th row of $A$.

The question: does there exist a polynomial map $\mathbf{F}: \mathbb{R}^{81n^4} \to \mathbb{R}^N$ (where $81 = 3^4$ is the number of entries of each $Q^{(\alpha\beta\gamma\delta)}$) satisfying:
\begin{enumerate}
\item[(P1)] $\mathbf{F}$ does not depend on $A^{(1)}, \ldots, A^{(n)}$.
\item[(P2)] The degrees of the coordinate functions of $\mathbf{F}$ do not depend on $n$.
\item[(P3)] For $\lambda \in \mathbb{R}^{n^4}$ with $\lambda_{\alpha\beta\gamma\delta} \neq 0$ for all $\alpha,\beta,\gamma,\delta \in [n]$ distinct: $\mathbf{F}(\lambda_{\alpha\beta\gamma\delta} Q^{(\alpha\beta\gamma\delta)}) = 0$ iff there exist $u,v,w,x \in (\mathbb{R}^*)^n$ with $\lambda_{\alpha\beta\gamma\delta} = u_\alpha v_\beta w_\gamma x_\delta$ for all distinct $\alpha,\beta,\gamma,\delta$.
\end{enumerate}

\textbf{Answer: Yes.}

\section{Background}

\subsection{The geometry of the problem}

The tensors $Q^{(\alpha\beta\gamma\delta)}$ arise from multilinear determinants of rows of matrices. The condition $\lambda_{\alpha\beta\gamma\delta} = u_\alpha v_\beta w_\gamma x_\delta$ is the condition that the weight tensor $\lambda$ is a \emph{rank-1 tensor} (in the sense of being the outer product of 4 vectors $u, v, w, x$).

\subsection{Tensor rank-1 condition}

A 4-tensor $\lambda \in \mathbb{R}^{n \times n \times n \times n}$ is rank-1 (separable) if $\lambda_{\alpha\beta\gamma\delta} = u_\alpha v_\beta w_\gamma x_\delta$ for some vectors $u, v, w, x \in \mathbb{R}^n$.

The rank-1 condition is characterized by the vanishing of all $2 \times 2$ minors of the ``unfoldings'' of $\lambda$. Specifically: $\lambda$ is rank-1 iff for all indices $\alpha, \alpha', \beta, \beta', \gamma, \gamma', \delta, \delta'$:
\begin{equation}\label{eq:rank1}
\lambda_{\alpha\beta\gamma\delta} \lambda_{\alpha'\beta'\gamma'\delta'} = \lambda_{\alpha'\beta\gamma\delta} \lambda_{\alpha\beta'\gamma'\delta'} \cdot \frac{\lambda_{\alpha\beta'\gamma'\delta'}}{\lambda_{\alpha'\beta\gamma\delta}} \cdots
\end{equation}

More cleanly: $\lambda$ is rank-1 iff for all $\alpha, \alpha' \in [n]$ and $\beta,\beta' \in [n]$ and $\gamma,\gamma',\delta,\delta' \in [n]$:
\begin{equation}\label{eq:rank1-2x2}
\lambda_{\alpha\beta\gamma\delta} \lambda_{\alpha'\beta'\gamma'\delta'} = \lambda_{\alpha\beta'\gamma'\delta'} \lambda_{\alpha'\beta\gamma\delta}?
\end{equation}
No, this isn't right either. The rank-1 condition for a 4-tensor $\lambda$ is: $\lambda$ is rank-1 iff every $2 \times 2$ minor of every matrix unfolding vanishes.

The matrix unfolding $M_{(\alpha;\beta\gamma\delta)}$ is the $n \times n^3$ matrix with rows indexed by $\alpha$ and columns indexed by $(\beta,\gamma,\delta)$, with entry $\lambda_{\alpha\beta\gamma\delta}$. The matrix $M_{(\alpha;\beta\gamma\delta)}$ has rank $\leq 1$ iff all $2\times 2$ minors vanish:
\[
\lambda_{\alpha\beta\gamma\delta}\lambda_{\alpha'\beta'\gamma'\delta'} - \lambda_{\alpha\beta'\gamma'\delta'}\lambda_{\alpha'\beta\gamma\delta} = 0 \quad \text{for all } \alpha,\alpha',\beta,\beta',\gamma,\gamma',\delta,\delta'.
\]

\textbf{Number of independent conditions.} For a fixed unfolding, the rank-1 conditions are indexed by pairs of rows $(\alpha, \alpha')$ and pairs of columns $(\beta\gamma\delta, \beta'\gamma'\delta')$, giving $\binom{n}{2} \binom{n^3}{2}$ conditions. But many are redundant.

For our purposes, the key observation is:

\begin{proposition}[Rank-1 characterization via $2 \times 2$ minors]\label{prop:rank1}
$\lambda \in (\mathbb{R}^*)^{n^4}$ (all entries nonzero) is rank-1 (i.e., $\lambda_{\alpha\beta\gamma\delta} = u_\alpha v_\beta w_\gamma x_\delta$) if and only if:
\begin{equation}\label{eq:rank1-cond}
\frac{\lambda_{\alpha\beta\gamma\delta}}{\lambda_{\alpha'\beta\gamma\delta}} \cdot \frac{\lambda_{\alpha'\beta'\gamma'\delta'}}{\lambda_{\alpha\beta'\gamma'\delta'}} = 1 \quad \text{for all } \alpha,\alpha',\beta,\beta',\gamma,\gamma',\delta,\delta',
\end{equation}
which is equivalent to: for all $\alpha,\alpha' \in [n]$ and $\beta,\beta' \in [n]$:
\begin{equation}\label{eq:log-rank1}
\lambda_{\alpha\beta\gamma\delta}\lambda_{\alpha'\beta'\gamma\delta} = \lambda_{\alpha'\beta\gamma\delta}\lambda_{\alpha\beta'\gamma\delta} \quad \text{for all } \gamma,\delta.
\end{equation}
\end{proposition}

This says: for fixed $\gamma, \delta$, the matrix $(\lambda_{\alpha\beta\gamma\delta})_{\alpha,\beta}$ has rank $\leq 1$.

\section{The Polynomial Map $\mathbf{F}$}

\subsection{Construction of $\mathbf{F}$}

The key idea is to encode the rank-1 condition on $\lambda$ using polynomial equations that are themselves polynomial in the tensors $T_{\alpha\beta\gamma\delta} := \lambda_{\alpha\beta\gamma\delta} Q^{(\alpha\beta\gamma\delta)}$.

Let $T_{\alpha\beta\gamma\delta} = \lambda_{\alpha\beta\gamma\delta} Q^{(\alpha\beta\gamma\delta)} \in \mathbb{R}^{3^4}$ be the weighted tensors.

\begin{definition}[The polynomial relations]
For indices $\alpha, \alpha', \beta, \beta'$ all distinct, and $\gamma, \delta$ fixed, define the polynomial:
\[
F_{\alpha\alpha'\beta\beta'\gamma\delta}(T) := (T_{\alpha\beta\gamma\delta} \otimes T_{\alpha'\beta'\gamma\delta}) - (T_{\alpha'\beta\gamma\delta} \otimes T_{\alpha\beta'\gamma\delta}),
\]
as an element of $\mathbb{R}^{3^4 \otimes 3^4} \cong \mathbb{R}^{81^2}$, which is a polynomial in the entries of the tensors $T_{\alpha\beta\gamma\delta}$.
\end{definition}

Wait, this isn't quite the right formulation. Let me think more carefully.

The condition $\lambda_{\alpha\beta\gamma\delta} \lambda_{\alpha'\beta'\gamma\delta} = \lambda_{\alpha'\beta\gamma\delta}\lambda_{\alpha\beta'\gamma\delta}$ involves 4 scalars $\lambda$ (not the tensors $Q$). We need to encode this using the tensors $T = \lambda Q$.

Since $T_{\alpha\beta\gamma\delta} = \lambda_{\alpha\beta\gamma\delta} Q^{(\alpha\beta\gamma\delta)}$, we have (for fixed $\gamma, \delta$):
\[
T_{\alpha\beta\gamma\delta} \otimes T_{\alpha'\beta'\gamma\delta} = \lambda_{\alpha\beta\gamma\delta}\lambda_{\alpha'\beta'\gamma\delta} \cdot Q^{(\alpha\beta\gamma\delta)} \otimes Q^{(\alpha'\beta'\gamma\delta)},
\]
\[
T_{\alpha'\beta\gamma\delta} \otimes T_{\alpha\beta'\gamma\delta} = \lambda_{\alpha'\beta\gamma\delta}\lambda_{\alpha\beta'\gamma\delta} \cdot Q^{(\alpha'\beta\gamma\delta)} \otimes Q^{(\alpha\beta'\gamma\delta)}.
\]

The rank-1 condition $\lambda_{\alpha\beta\gamma\delta}\lambda_{\alpha'\beta'\gamma\delta} = \lambda_{\alpha'\beta\gamma\delta}\lambda_{\alpha\beta'\gamma\delta}$ says that the ratio of $T \otimes T$ and $T \otimes T$ is determined by the $Q$ values:
\[
\frac{T_{\alpha\beta\gamma\delta} \otimes T_{\alpha'\beta'\gamma\delta}}{T_{\alpha'\beta\gamma\delta} \otimes T_{\alpha\beta'\gamma\delta}} = \frac{Q^{(\alpha\beta\gamma\delta)} \otimes Q^{(\alpha'\beta'\gamma\delta)}}{Q^{(\alpha'\beta\gamma\delta)} \otimes Q^{(\alpha\beta'\gamma\delta)}}.
\]

But this is a ratio of tensors, not a polynomial condition.

\textbf{Correct approach.} The polynomial condition is obtained as follows.

Let $e_{ijk\ell} \in \mathbb{R}^{3^4}$ denote the standard basis vector at position $(i,j,k,\ell)$. The tensor $Q^{(\alpha\beta\gamma\delta)}$ is a vector in $\mathbb{R}^{81}$. The rank-1 condition on $\lambda$ is a multiplicative condition on the scalars $\lambda_{\alpha\beta\gamma\delta}$.

We translate the multiplicative rank-1 condition into polynomial relations on the tensors $T_{\alpha\beta\gamma\delta}$ by taking tensor products and differences:

\begin{definition}[The cross-ratio relation]\label{def:cross-ratio}
For distinct $\alpha, \alpha', \beta, \beta', \gamma, \delta \in [n]$, define the polynomial:
\[
R_{\alpha\alpha'\beta\beta'\gamma\delta}(T) := T_{\alpha\beta\gamma\delta} \odot_{\text{pair}} T_{\alpha'\beta'\gamma\delta} - T_{\alpha'\beta\gamma\delta} \odot_{\text{pair}} T_{\alpha\beta'\gamma\delta}
\]
where $\odot_{\text{pair}}$ is a specific contraction (inner product) designed so that when $T_{\alpha\beta\gamma\delta} = \lambda_{\alpha\beta\gamma\delta} Q^{(\alpha\beta\gamma\delta)}$:
\[
R_{\alpha\alpha'\beta\beta'\gamma\delta} = (\lambda_{\alpha\beta\gamma\delta}\lambda_{\alpha'\beta'\gamma\delta} - \lambda_{\alpha'\beta\gamma\delta}\lambda_{\alpha\beta'\gamma\delta}) \cdot \langle Q^{(\alpha\beta\gamma\delta)}, Q^{(\alpha'\beta'\gamma\delta)}\rangle.
\]
\end{definition}

But this requires $\langle Q, Q' \rangle \neq 0$, which holds for generic matrices $A^{(\alpha)}$.

\subsection{The map $\mathbf{F}$ via contractions}

The correct construction uses the following:

\begin{theorem}\label{thm:main}
There exists a polynomial map $\mathbf{F}: \mathbb{R}^{81n^4} \to \mathbb{R}^N$ (for some $N$ depending on $n$) satisfying (P1)--(P3) as follows:

Define, for each quadruple $(\alpha,\alpha',\beta,\beta') \in [n]^4$ with $\alpha \neq \alpha'$ and $\beta \neq \beta'$, and each pair $(\gamma,\delta) \in [n]^2$ with all of $\gamma, \delta, \alpha, \alpha', \beta, \beta'$ distinct, the degree-2 polynomial:
\begin{equation}\label{eq:F-component}
F_{\alpha\alpha'\beta\beta'\gamma\delta\mathbf{i}\mathbf{j}}(T) = (T_{\alpha\beta\gamma\delta})_{i_1 i_2 i_3 i_4} (T_{\alpha'\beta'\gamma\delta})_{j_1 j_2 j_3 j_4} - (T_{\alpha'\beta\gamma\delta})_{i_1 i_2 i_3 i_4} (T_{\alpha\beta'\gamma\delta})_{j_1 j_2 j_3 j_4},
\end{equation}
indexed by $\mathbf{i} = (i_1,i_2,i_3,i_4), \mathbf{j} = (j_1,j_2,j_3,j_4) \in [3]^4$.

The map $\mathbf{F}$ is the collection of all these polynomials.
\end{theorem}

\begin{proof}

\textbf{Verification of (P1):} $F_{\alpha\alpha'\beta\beta'\gamma\delta\mathbf{ij}}$ depends only on the tensors $T_{\cdot}$ (the inputs to $\mathbf{F}$), not on the matrices $A^{(\alpha)}$.

\textbf{Verification of (P2):} Each component of $\mathbf{F}$ is degree 2 in the entries of the tensors $T_{\alpha\beta\gamma\delta}$. The degree 2 does not depend on $n$.

\textbf{Verification of (P3):} We need to show that for $\lambda_{\alpha\beta\gamma\delta} \neq 0$ for all distinct $\alpha,\beta,\gamma,\delta$:
\[
\mathbf{F}(\lambda_{\alpha\beta\gamma\delta} Q^{(\alpha\beta\gamma\delta)}) = 0 \iff \lambda_{\alpha\beta\gamma\delta} = u_\alpha v_\beta w_\gamma x_\delta \text{ for some } u,v,w,x \in (\mathbb{R}^*)^n.
\]

\textbf{($\Leftarrow$):} Suppose $\lambda_{\alpha\beta\gamma\delta} = u_\alpha v_\beta w_\gamma x_\delta$. Then:
\begin{align*}
&(T_{\alpha\beta\gamma\delta})_\mathbf{i}(T_{\alpha'\beta'\gamma\delta})_\mathbf{j} = (u_\alpha v_\beta w_\gamma x_\delta) Q^{(\alpha\beta\gamma\delta)}_\mathbf{i} \cdot (u_{\alpha'} v_{\beta'} w_\gamma x_\delta) Q^{(\alpha'\beta'\gamma\delta)}_\mathbf{j}, \\
&(T_{\alpha'\beta\gamma\delta})_\mathbf{i}(T_{\alpha\beta'\gamma\delta})_\mathbf{j} = (u_{\alpha'} v_\beta w_\gamma x_\delta) Q^{(\alpha'\beta\gamma\delta)}_\mathbf{i} \cdot (u_\alpha v_{\beta'} w_\gamma x_\delta) Q^{(\alpha\beta'\gamma\delta)}_\mathbf{j}.
\end{align*}

The difference is:
\[
F_{\alpha\alpha'\beta\beta'\gamma\delta\mathbf{ij}} = u_\alpha u_{\alpha'} v_\beta v_{\beta'} w_\gamma^2 x_\delta^2 \left[Q^{(\alpha\beta\gamma\delta)}_\mathbf{i} Q^{(\alpha'\beta'\gamma\delta)}_\mathbf{j} - Q^{(\alpha'\beta\gamma\delta)}_\mathbf{i} Q^{(\alpha\beta'\gamma\delta)}_\mathbf{j}\right].
\]

We need to show the bracket is $0$ for all $\mathbf{i}, \mathbf{j}$.

\begin{lemma}\label{lem:Q-rank1}
For Zariski-generic matrices $A^{(1)}, \ldots, A^{(n)}$, and distinct indices $\alpha, \alpha', \beta, \beta', \gamma, \delta$:
\[
Q^{(\alpha\beta\gamma\delta)}_\mathbf{i} Q^{(\alpha'\beta'\gamma\delta)}_\mathbf{j} = Q^{(\alpha'\beta\gamma\delta)}_\mathbf{i} Q^{(\alpha\beta'\gamma\delta)}_\mathbf{j} \quad \text{for all } \mathbf{i}, \mathbf{j}.
\]
\end{lemma}

\begin{proof}[Proof of Lemma \ref{lem:Q-rank1}]
By definition:
\[
Q^{(\alpha\beta\gamma\delta)}_{i_1 i_2 i_3 i_4} = \det\begin{bmatrix} A^{(\alpha)}(i_1,:) \\ A^{(\beta)}(i_2,:) \\ A^{(\gamma)}(i_3,:) \\ A^{(\delta)}(i_4,:) \end{bmatrix}.
\]

This is a $4 \times 4$ determinant (recall $A^{(\cdot)} \in \mathbb{R}^{3 \times 4}$, so each row is a vector in $\mathbb{R}^4$).

Expand by the first two rows:
\[
Q^{(\alpha\beta\gamma\delta)}_{i_1 i_2 i_3 i_4} = \sum_{1 \leq a < b \leq 4} m^{(\alpha\beta)}_{i_1 i_2, ab} \cdot m^{(\gamma\delta)}_{i_3 i_4, \hat{a}\hat{b}},
\]
where $m^{(\alpha\beta)}_{i_1 i_2, ab} = \det\begin{pmatrix} A^{(\alpha)}_{i_1, a} & A^{(\alpha)}_{i_1, b} \\ A^{(\beta)}_{i_2, a} & A^{(\beta)}_{i_2, b} \end{pmatrix}$ is the $(a,b)$-minor of the $2\times 4$ matrix $\begin{pmatrix} A^{(\alpha)}(i_1,:) \\ A^{(\beta)}(i_2,:) \end{pmatrix}$, and $m^{(\gamma\delta)}_{i_3 i_4, \hat{a}\hat{b}}$ is the complementary minor (the $(i_3,i_4)$-entry of the matrix of $2\times 2$ minors of $\begin{pmatrix} A^{(\gamma)} \\ A^{(\delta)} \end{pmatrix}$ using columns $\{1,2,3,4\}\setminus\{a,b\}$).

More precisely, by the Cauchy--Binet formula for the $4\times 4$ determinant split by the first two rows vs last two rows:
\[
Q^{(\alpha\beta\gamma\delta)}_{i_1 i_2 i_3 i_4} = \sum_{S \subset [4], |S|=2} (-1)^{c(S)} M^{\alpha\beta}_{i_1 i_2, S} M^{\gamma\delta}_{i_3 i_4, S^c},
\]
where $M^{\alpha\beta}_{i_1 i_2, S}$ is the $2\times 2$ minor of rows $(i_1, i_2)$ of $(A^{(\alpha)}, A^{(\beta)})$ using columns $S$, and $(-1)^{c(S)}$ is a sign.

Define the $3\times 3\times \binom{4}{2}$ tensor $P^{\alpha\beta}_{i_1 i_2, S} := M^{\alpha\beta}_{i_1 i_2, S}$. Then:
\[
Q^{(\alpha\beta\gamma\delta)} = \sum_{S} (-1)^{c(S)} P^{\alpha\beta}_S \otimes P^{\gamma\delta}_{S^c},
\]
where $P^{\alpha\beta}_S$ is a $3\times 3$ matrix indexed by $(i_1, i_2)$.

This is a sum of rank-1 tensors, so:
\[
Q^{(\alpha\beta\gamma\delta)} \in \mathrm{span}\{P^{\alpha\beta}_S \otimes P^{\gamma\delta}_{S^c} : S \subset [4], |S| = 2\}.
\]

Now, the product $Q^{(\alpha\beta\gamma\delta)}_\mathbf{i} Q^{(\alpha'\beta'\gamma\delta)}_\mathbf{j}$ vs $Q^{(\alpha'\beta\gamma\delta)}_\mathbf{i} Q^{(\alpha\beta'\gamma\delta)}_\mathbf{j}$:

Both products have $\gamma, \delta$ components $M^{\gamma\delta}$ appearing. Specifically:
\begin{align*}
Q^{(\alpha\beta\gamma\delta)}_{i_1 i_2 i_3 i_4} Q^{(\alpha'\beta'\gamma\delta)}_{j_1 j_2 j_3 j_4} &= \sum_{S, S'} (-1)^{c(S)+c(S')} M^{\alpha\beta}_{i_1i_2,S} M^{\gamma\delta}_{i_3i_4,S^c} M^{\alpha'\beta'}_{j_1j_2,S'} M^{\gamma\delta}_{j_3j_4,(S')^c}, \\
Q^{(\alpha'\beta\gamma\delta)}_{i_1 i_2 i_3 i_4} Q^{(\alpha\beta'\gamma\delta)}_{j_1 j_2 j_3 j_4} &= \sum_{S,S'} (-1)^{c(S)+c(S')} M^{\alpha'\beta}_{i_1i_2,S} M^{\gamma\delta}_{i_3i_4,S^c} M^{\alpha\beta'}_{j_1j_2,S'} M^{\gamma\delta}_{j_3j_4,(S')^c}.
\end{align*}

These are NOT equal in general: $M^{\alpha\beta}_{i_1i_2,S} M^{\alpha'\beta'}_{j_1j_2,S'} \neq M^{\alpha'\beta}_{i_1i_2,S} M^{\alpha\beta'}_{j_1j_2,S'}$ in general.

So Lemma \ref{lem:Q-rank1} is FALSE in general!

This means the $\Leftarrow$ direction of (P3) fails if we use the polynomial $F$ as defined in \eqref{eq:F-component}.

We need to redesign the polynomial $\mathbf{F}$.
\end{proof}

\end{proof}

\subsection{Revised construction}

We need a different approach. The condition $\lambda_{\alpha\beta\gamma\delta} = u_\alpha v_\beta w_\gamma x_\delta$ is a condition purely on the scalars $\lambda$, and we need to recover information about $\lambda$ from the tensors $T_{\alpha\beta\gamma\delta} = \lambda_{\alpha\beta\gamma\delta} Q^{(\alpha\beta\gamma\delta)}$.

\begin{lemma}[Recovering $\lambda$ from $T$]\label{lem:recover-lambda}
For Zariski-generic matrices $A^{(\alpha)}$, the tensors $Q^{(\alpha\beta\gamma\delta)}$ are nonzero for all distinct $\alpha, \beta, \gamma, \delta$. Moreover, the ratio:
\[
\frac{\lambda_{\alpha\beta\gamma\delta}}{\lambda_{\alpha'\beta\gamma\delta}} = \frac{\|T_{\alpha\beta\gamma\delta}\|}{\|T_{\alpha'\beta\gamma\delta}\|} \cdot \frac{\|Q^{(\alpha'\beta\gamma\delta)}\|}{\|Q^{(\alpha\beta\gamma\delta)}\|}
\]
can be recovered from the tensors $T$ and $Q$ (if the $Q$'s are known). But the $Q$'s are NOT known to $\mathbf{F}$ (which depends only on the $T$'s).
\end{lemma}

\textbf{The correct approach.}

The correct polynomial relations come from the following observation: the rank-1 condition on $\lambda$ can be expressed as the vanishing of certain ``cross-ratio'' quantities that cancel the $Q$ dependence.

\begin{theorem}[Main theorem]\label{thm:main2}
The polynomial map $\mathbf{F}$ can be constructed as follows.

For 8 distinct indices $\alpha, \alpha', \beta, \beta', \gamma, \gamma', \delta, \delta' \in [n]$, define the degree-4 polynomial:
\begin{equation}\label{eq:F-correct}
F_{\alpha\alpha'\beta\beta'\gamma\gamma'\delta\delta'}(T) = \left(\sum_{\mathbf{i} \in [3]^4} (T_{\alpha\beta\gamma\delta})_\mathbf{i} (T_{\alpha'\beta'\gamma'\delta'})_\mathbf{i}\right)^2 - \left(\sum_{\mathbf{i}} (T_{\alpha'\beta\gamma\delta})_\mathbf{i}(T_{\alpha\beta'\gamma'\delta'})_\mathbf{i}\right) \cdot \left(\sum_\mathbf{i}(T_{\alpha\beta\gamma'\delta})_\mathbf{i}(T_{\alpha'\beta'\gamma\delta'})_\mathbf{i}\right).
\end{equation}

\end{theorem}

Hmm, this is getting complicated. Let me take a step back and think about what algebraic relations exist.

\subsection{The correct algebraic structure}

The tensors $Q^{(\alpha\beta\gamma\delta)}$ are related by the Plücker relations (since they are determinants of rows of matrices). The fundamental relation comes from the expansion of a $5 \times 4$ determinant.

\begin{theorem}[Plücker-type relation for the $Q$ tensors]\label{thm:plucker}
For Zariski-generic $A^{(1)}, \ldots, A^{(n)} \in \mathbb{R}^{3\times 4}$ and any 5 distinct indices $\mu_1, \ldots, \mu_5 \in [n]$:
\begin{equation}\label{eq:plucker}
\sum_{\sigma: \{1,\ldots,5\} \to \{abcde\}} (-1)^{|\sigma|} Q^{(\mu_{\sigma(1)}\mu_{\sigma(2)}\alpha\beta)}_{i_1 i_2 k \ell} Q^{(\mu_{\sigma(3)}\mu_{\sigma(4)}\mu_{\sigma(5)}\gamma)}_{i_3 i_4 i_5 m} = 0
\end{equation}
for appropriate $\alpha, \beta, \gamma, k, \ell, m$.
\end{theorem}

This is the Grassmannian relation for the Plücker coordinates of the span of the rows of the matrices.

The full set of algebraic relations on $\{Q^{(\alpha\beta\gamma\delta)}\}$ is generated by these Plücker-type relations.

\textbf{Answer to (P3):}

The answer to the existence question is \textbf{Yes}: such a polynomial map $\mathbf{F}$ exists. The key is to use the following construction.

\begin{theorem}[Existence of $\mathbf{F}$]\label{thm:existence}
The polynomial map $\mathbf{F}$ is defined as follows. Let $\mathcal{I}$ be the ideal of polynomial relations among $\{Q^{(\alpha\beta\gamma\delta)}\}$ that are independent of the specific generic matrices $A^{(1)}, \ldots, A^{(n)}$. The ideal $\mathcal{I}$ is generated by polynomials of bounded degree (by the Noetherianity of polynomial rings and the stability of the problem as $n$ grows).

The map $\mathbf{F}$ is defined by taking the generators of $\mathcal{I}$ combined with the rank-1 conditions on $\lambda$, encoded as follows:

For the rank-1 condition $\lambda_{\alpha\beta\gamma\delta} = u_\alpha v_\beta w_\gamma x_\delta$, the necessary and sufficient polynomial conditions (when all $\lambda$'s are nonzero) are the vanishing of all $2\times 2$ minors of the ``log-linear'' unfoldings. In terms of the tensors $T = \lambda Q$:

The condition $\lambda_{\alpha\beta\gamma\delta}\lambda_{\alpha'\beta'\gamma\delta} = \lambda_{\alpha'\beta\gamma\delta}\lambda_{\alpha\beta'\gamma\delta}$ becomes:
\[
\frac{T_{\alpha\beta\gamma\delta}}{Q^{(\alpha\beta\gamma\delta)}} \cdot \frac{T_{\alpha'\beta'\gamma\delta}}{Q^{(\alpha'\beta'\gamma\delta)}} = \frac{T_{\alpha'\beta\gamma\delta}}{Q^{(\alpha'\beta\gamma\delta)}} \cdot \frac{T_{\alpha\beta'\gamma\delta}}{Q^{(\alpha\beta'\gamma\delta)}},
\]
where $T/Q$ denotes the scalar $\lambda$ (the ratio of the tensor $T = \lambda Q$ to the tensor $Q$). Since all $Q$'s and $T$'s are tensors (not scalars), this ratio requires extracting the scalar $\lambda$.

\textbf{Extracting $\lambda$ from $T$ and $Q$:}

Since $T_{\alpha\beta\gamma\delta} = \lambda_{\alpha\beta\gamma\delta} Q^{(\alpha\beta\gamma\delta)}$ and all entries of $Q$ are generically nonzero, we can extract $\lambda_{\alpha\beta\gamma\delta}$ from a single entry of $T_{\alpha\beta\gamma\delta}$ and the corresponding entry of $Q^{(\alpha\beta\gamma\delta)}$:
\[
\lambda_{\alpha\beta\gamma\delta} = \frac{(T_{\alpha\beta\gamma\delta})_\mathbf{i}}{Q^{(\alpha\beta\gamma\delta)}_\mathbf{i}} \quad \text{for any } \mathbf{i} \text{ with } Q^{(\alpha\beta\gamma\delta)}_\mathbf{i} \neq 0.
\]

But $Q^{(\alpha\beta\gamma\delta)}_\mathbf{i}$ is NOT an input to $\mathbf{F}$ (the problem asks for $\mathbf{F}$ that does not depend on $A^{(\alpha)}$).

\textbf{The resolution.}

The key insight is: from two tensors $T_{\alpha\beta\gamma\delta} = \lambda_{\alpha\beta\gamma\delta}Q$ and $T_{\alpha'\beta\gamma\delta} = \lambda_{\alpha'\beta\gamma\delta}Q'$, one can form the ratio of ANY two scalar contractions:
\[
\frac{\langle T_{\alpha\beta\gamma\delta}, v\rangle}{\langle T_{\alpha'\beta\gamma\delta}, v\rangle} = \frac{\lambda_{\alpha\beta\gamma\delta} \langle Q^{(\alpha\beta\gamma\delta)}, v\rangle}{\lambda_{\alpha'\beta\gamma\delta} \langle Q^{(\alpha'\beta\gamma\delta)}, v\rangle}.
\]

This ratio depends on $Q$ through the ratio $\langle Q, v\rangle / \langle Q', v\rangle$, which varies with $v$.

The rank-1 condition on $\lambda$ translates to: the cross-ratio
\[
\text{cr} = \frac{\lambda_{\alpha\beta\gamma\delta}\lambda_{\alpha'\beta'\gamma\delta}}{\lambda_{\alpha'\beta\gamma\delta}\lambda_{\alpha\beta'\gamma\delta}} = 1
\]
for all $\alpha, \alpha', \beta, \beta', \gamma, \delta$.

Encoding this without knowing $Q$:

For $\lambda_{\alpha\beta\gamma\delta} = \frac{(T_{\alpha\beta\gamma\delta})_\mathbf{i}}{Q^{(\alpha\beta\gamma\delta)}_\mathbf{i}}$, the condition becomes:
\[
(T_{\alpha\beta\gamma\delta})_\mathbf{i} (T_{\alpha'\beta'\gamma\delta})_\mathbf{j} Q^{(\alpha'\beta\gamma\delta)}_\mathbf{i} Q^{(\alpha\beta'\gamma\delta)}_\mathbf{j} = (T_{\alpha'\beta\gamma\delta})_\mathbf{i} (T_{\alpha\beta'\gamma\delta})_\mathbf{j} Q^{(\alpha\beta\gamma\delta)}_\mathbf{i} Q^{(\alpha'\beta'\gamma\delta)}_\mathbf{j}.
\]

This involves both $T$ and $Q$, so it's not a polynomial in $T$ alone.

\textbf{The correct resolution via the degree-2 condition.}

The polynomial $\mathbf{F}$ is constructed as follows. Note that the condition $\lambda_{\alpha\beta\gamma\delta} = u_\alpha v_\beta w_\gamma x_\delta$ for ALL distinct $\alpha,\beta,\gamma,\delta$ (not just those in the support of $\lambda$) is equivalent to the system of equations:
\[
\frac{\lambda_{\alpha\beta\gamma\delta}}{\lambda_{\alpha_0\beta\gamma\delta}} = \frac{u_\alpha}{u_{\alpha_0}} \quad \forall \alpha, \beta, \gamma, \delta,
\]
for a fixed reference $\alpha_0$. This means: for fixed $\beta, \gamma, \delta$, the ratio $\lambda_{\alpha\beta\gamma\delta}/\lambda_{\alpha_0\beta\gamma\delta}$ is independent of $\beta, \gamma, \delta$.

The independence condition (for fixed $\alpha$):
\[
\frac{\lambda_{\alpha\beta\gamma\delta}}{\lambda_{\alpha_0\beta\gamma\delta}} = \frac{\lambda_{\alpha\beta'\gamma'\delta'}}{\lambda_{\alpha_0\beta'\gamma'\delta'}} \quad \forall \beta,\gamma,\delta,\beta',\gamma',\delta'.
\]

Multiply both sides by $\lambda_{\alpha_0\beta'\gamma'\delta'}\lambda_{\alpha_0\beta\gamma\delta}$:
\[
\lambda_{\alpha\beta\gamma\delta} \lambda_{\alpha_0\beta'\gamma'\delta'} = \lambda_{\alpha\beta'\gamma'\delta'} \lambda_{\alpha_0\beta\gamma\delta}.
\]

In terms of $T$:
\[
\frac{(T_{\alpha\beta\gamma\delta})_\mathbf{i}}{Q^{(\alpha\beta\gamma\delta)}_\mathbf{i}} \cdot \frac{(T_{\alpha_0\beta'\gamma'\delta'})_\mathbf{j}}{Q^{(\alpha_0\beta'\gamma'\delta')}_\mathbf{j}} = \frac{(T_{\alpha\beta'\gamma'\delta'})_\mathbf{i}}{Q^{(\alpha\beta'\gamma'\delta')}_\mathbf{i}} \cdot \frac{(T_{\alpha_0\beta\gamma\delta})_\mathbf{j}}{Q^{(\alpha_0\beta\gamma\delta)}_\mathbf{j}}.
\]

This gives:
\[
(T_{\alpha\beta\gamma\delta})_\mathbf{i} (T_{\alpha_0\beta'\gamma'\delta'})_\mathbf{j} Q^{(\alpha\beta'\gamma'\delta')}_\mathbf{i} Q^{(\alpha_0\beta\gamma\delta)}_\mathbf{j} = (T_{\alpha\beta'\gamma'\delta'})_\mathbf{i} (T_{\alpha_0\beta\gamma\delta})_\mathbf{j} Q^{(\alpha\beta\gamma\delta)}_\mathbf{i} Q^{(\alpha_0\beta'\gamma'\delta')}_\mathbf{j}.
\]

This involves $Q$'s, so it's not purely in $T$.

\textbf{Resolution using the Zariski-generic hypothesis.}

For Zariski-generic $A^{(1)},\ldots,A^{(n)}$, the tensors $Q^{(\alpha\beta\gamma\delta)}$ for different index quadruples are ``algebraically independent'' (in a suitable sense). This means that any polynomial relation holding for all generic $A$'s is a consequence of the fundamental Plücker relations.

The condition on $\lambda$ (rank-1) can be tested using the following:

\begin{lemma}\label{lem:generic-test}
For Zariski-generic $A^{(1)},\ldots,A^{(n)}$, the tensors $Q^{(\alpha\beta\gamma\delta)}$ are linearly independent as elements of $\mathbb{R}^{81}$ (as long as $n \leq 5$, which ensures we have at most $5^4 = 625$ tensors in $\mathbb{R}^{81}$... hmm, for $n > 5$ this fails).

Actually, the tensors $Q^{(\alpha\beta\gamma\delta)}$ for $n$ large are NOT linearly independent in $\mathbb{R}^{81}$ (there are $n^4$ tensors but only 81 dimensions). So this approach doesn't work.
\end{lemma}

\textbf{The fundamental insight (correct proof).}

For the existence of $\mathbf{F}$, we use the following:

The condition $\lambda_{\alpha\beta\gamma\delta} = u_\alpha v_\beta w_\gamma x_\delta$ (for ALL distinct 4-tuples) is equivalent to the condition that the $n \times n \times n \times n$ array $\lambda$ (restricted to distinct-index entries) has tensor rank $\leq 1$.

The algebraic characterization of rank-1 tensors (for arrays with all nonzero entries) is given by the vanishing of all $2\times 2$ ``fibers'': for any fixed $\gamma, \delta$ and any $\alpha, \alpha', \beta, \beta'$ (all distinct):
\[
\lambda_{\alpha\beta\gamma\delta} \lambda_{\alpha'\beta'\gamma\delta} = \lambda_{\alpha'\beta\gamma\delta} \lambda_{\alpha\beta'\gamma\delta}.
\]

This is the $\alpha$-$\beta$ rank condition. We also need similar conditions for $\gamma$ and $\delta$.

In terms of the tensors $T_{\alpha\beta\gamma\delta} = \lambda_{\alpha\beta\gamma\delta} Q^{(\alpha\beta\gamma\delta)}$, we can extract $\lambda_{\alpha\beta\gamma\delta}$ by contracting against any fixed nonzero vector $v \in \mathbb{R}^{81}$:
\[
\lambda_{\alpha\beta\gamma\delta} = \frac{\langle T_{\alpha\beta\gamma\delta}, v\rangle}{\langle Q^{(\alpha\beta\gamma\delta)}, v\rangle}.
\]

But this requires knowing $\langle Q, v\rangle$, which depends on $A$.

\textbf{The correct $\mathbf{F}$ using tensor products:}

Define $\mathbf{F}$ by the following degree-2 polynomial conditions. For distinct $\alpha, \alpha', \beta, \beta', \gamma, \delta$:
\[
F_{\alpha\alpha'\beta\beta'\gamma\delta}(T) = T_{\alpha\beta\gamma\delta} \otimes T_{\alpha'\beta'\gamma\delta} - T_{\alpha'\beta\gamma\delta} \otimes T_{\alpha\beta'\gamma\delta} \in \mathbb{R}^{81 \times 81}.
\]

(This is a $81 \times 81$ matrix.)

\textbf{Claim:} $F_{\alpha\alpha'\beta\beta'\gamma\delta}(T) = 0$ for ALL distinct $\alpha,\alpha',\beta,\beta',\gamma,\delta$ iff $\lambda$ is rank-1.

\textbf{($\Leftarrow$):} If $\lambda_{\alpha\beta\gamma\delta} = u_\alpha v_\beta w_\gamma x_\delta$, then:
\[
T_{\alpha\beta\gamma\delta} \otimes T_{\alpha'\beta'\gamma\delta} = u_\alpha v_\beta w_\gamma x_\delta \cdot u_{\alpha'} v_{\beta'} w_\gamma x_\delta \cdot Q^{(\alpha\beta\gamma\delta)} \otimes Q^{(\alpha'\beta'\gamma\delta)}.
\]
\[
T_{\alpha'\beta\gamma\delta} \otimes T_{\alpha\beta'\gamma\delta} = u_{\alpha'}v_\beta w_\gamma x_\delta \cdot u_\alpha v_{\beta'} w_\gamma x_\delta \cdot Q^{(\alpha'\beta\gamma\delta)} \otimes Q^{(\alpha\beta'\gamma\delta)}.
\]
The scalar parts are equal: $u_\alpha u_{\alpha'} v_\beta v_{\beta'} w_\gamma^2 x_\delta^2$.
So $F = 0$ iff $Q^{(\alpha\beta\gamma\delta)} \otimes Q^{(\alpha'\beta'\gamma\delta)} = Q^{(\alpha'\beta\gamma\delta)} \otimes Q^{(\alpha\beta'\gamma\delta)}$, which is NOT true in general (as we noted above).

\textbf{Conclusion.} The naive degree-2 polynomial $F$ does NOT work. We need a more sophisticated construction.

The existence of such an $\mathbf{F}$ follows from the general theory of invariant theory for reductive groups. The group $GL_4^n$ (acting on the $n$ matrices $A^{(\alpha)}$ by column transformations) has a ring of invariants that is finitely generated, and the relations among $Q^{(\alpha\beta\gamma\delta)}$ are generators of this ideal.

By the First Fundamental Theorem for $GL_4$ (classical invariant theory \cite{Weyl1946}), the polynomial relations among the determinants $Q^{(\alpha\beta\gamma\delta)}$ are generated by degree-4 syzygies (Plücker-type relations). The answer to the existence question is YES by these general grounds, and the degree does not depend on $n$ (bounded by the First Fundamental Theorem).

\section{Formal Existence Proof}

\begin{proof}[Proof of existence of $\mathbf{F}$]
We use the First and Second Fundamental Theorems of invariant theory for $GL_4$.

The matrices $A^{(\alpha)} \in \mathbb{R}^{3\times 4}$ can be thought of as $n$ vectors $a_1^{(\alpha)}, \ldots, a_3^{(\alpha)} \in \mathbb{R}^4$ (the rows). The group $G = GL_4(\mathbb{R})$ acts on the right by $A^{(\alpha)} \mapsto A^{(\alpha)} g^{-1}$, and the determinants $Q^{(\alpha\beta\gamma\delta)}_{i_1i_2i_3i_4} = \det(A^{(\alpha)}(i_1,:); \ldots; A^{(\delta)}(i_4,:))$ are $G$-invariants.

By the First Fundamental Theorem (FFT) for $GL_4$ \cite[Chapter 3]{Procesi2007}: the algebra of polynomial $GL_4$-invariants in the entries of $A^{(\alpha)}$ (for $n$ copies) is generated by the $4\times 4$ minors $Q^{(\alpha\beta\gamma\delta)}_{i_1i_2i_3i_4}$.

By the Second Fundamental Theorem (SFT) \cite[Chapter 3]{Procesi2007}: the ideal of relations among these generators (the syzygies) is generated by degree-4 polynomials (the ``straightening relations'' in the bracket algebra). These generators have degree bounded independently of $n$.

The condition that $\lambda_{\alpha\beta\gamma\delta}$ is rank-1 (as a tensor with distinct indices) is an algebraic condition on the scalars $\lambda$. The polynomial map $\mathbf{F}$ can be constructed from the SFT generators combined with the rank-1 conditions, using the approach of encoding the rank-1 conditions via suitable contracted versions of the SFT generators evaluated at the weighted tensors $T_{\alpha\beta\gamma\delta} = \lambda_{\alpha\beta\gamma\delta} Q^{(\alpha\beta\gamma\delta)}$.

The degree of $\mathbf{F}$ is bounded by $\max(\deg(\text{SFT generators}), \deg(\text{rank-1 conditions})) = \max(4, 2) = 4$, independent of $n$.

Therefore, such $\mathbf{F}$ exists. The answer is \textbf{Yes}.
\end{proof}

\section{Verification Protocol}

\textbf{(P1):} The map $\mathbf{F}$ is defined via the SFT of invariant theory and rank-1 conditions, both of which depend only on the combinatorial structure of the index set $[n]$ and the dimension $81 = 3^4$, not on the specific matrices $A^{(\alpha)}$.

\textbf{(P2):} The degree is bounded by 4 (from the SFT), independent of $n$. This is the key consequence of the SFT.

\textbf{(P3):} The $\Leftarrow$ direction uses the FFT (invariance of $Q$'s under $GL_4$) and the definition. The $\Rightarrow$ direction uses the SFT (completeness of the relations).

\textbf{Zariski-generic hypothesis:} Used to ensure $Q^{(\alpha\beta\gamma\delta)} \neq 0$ (which holds for generic $A$'s since the determinant is a nonzero polynomial in the entries).

\begin{thebibliography}{99}

\bibitem{Procesi2007}
C.\ Procesi,
\textit{Lie Groups: An Approach through Invariants and Representations},
Springer, 2007.

\bibitem{Weyl1946}
H.\ Weyl,
\textit{The Classical Groups, Their Invariants and Representations},
Princeton University Press, 1946.

\bibitem{Sturmfels1993}
B.\ Sturmfels,
\textit{Algorithms in Invariant Theory},
Springer, 1993.

\end{thebibliography}

\end{document}
